\documentclass[11pt]{article}
\usepackage[margin=1in]{geometry}
\usepackage[T1]{fontenc}\usepackage[utf8]{inputenc}
\usepackage{lmodern,amsmath,amssymb,amsthm,mathtools,booktabs,microtype}
\usepackage[colorlinks=true,linkcolor=blue,citecolor=blue,urlcolor=blue]{hyperref}
\hypersetup{pdftitle={Agnostic learning with layerwise heavy-tailed neural-network priors},pdfauthor={ChatGPT-6 Astra; Claude Opus 5.5; Kaarel Hänni; Dmitry Vaintrob}}
\newtheorem{theorem}{Theorem}[section]\newtheorem{lemma}[theorem]{Lemma}
\newtheorem{corollary}[theorem]{Corollary}\newtheorem{proposition}[theorem]{Proposition}
\theoremstyle{remark}\newtheorem{remark}[theorem]{Remark}
\newcommand{\E}{\mathbb E}\newcommand{\R}{\mathbb R}\newcommand{\Pp}{\mathbb P}
\newcommand{\KL}{\operatorname{KL}}\newcommand{\clip}{\operatorname{clip}}
\newcommand{\whR}{\widehat R}\newcommand{\calG}{\mathcal G}
\title{Agnostic learning with layerwise heavy-tailed\\neural-network priors}
\author{ChatGPT-6 Astra \and Claude Opus 5.5 \and Kaarel H\"anni \and
Dmitry Vaintrob\thanks{Authors are listed in alphabetical order.}}
\date{17 November 2026}
\begin{document}\maketitle
\begin{abstract}
We prove a common agnostic learning theorem for dense Bayesian neural networks with any fixed globally Lipschitz hidden activation, including ReLU, GELU, tanh, sigmoid and softplus. The prior has independent continuous heavy-tailed coordinates in one fixed architecture. Tail indices decrease strictly through the network; a row of fan-in $F$ and tail index $\alpha$ uses scale $F^{-1/\alpha}$. A fixed bounded output link, such as tanh or clipping, enters a Gaussian likelihood whose variance is independent of sample size and learner width. At every finite learner width above comparator width, the ordinary final posterior mean achieves the label-noise floor plus $(1+\eta)$ times the best smaller network's approximation error, with a polynomial remainder. If the second-layer tail index is $a<1$, the remainder is of order $n^{-(2-a)/2}$ for every globally Lipschitz activation, with no smoothness requirement. A complementary all-data expectation theorem has coefficient exactly one on comparator risk, with an explicit approximation-dependent rate. The comparator-complexity factor is polynomial in width, coefficient bound and input dimension at fixed depth. The proof conditionally truncates background weights according to preceding feature magnitudes. Decreasing tail indices make the required fractional moments finite, while the multiplicity of small-network embeddings compensates their coordinate probabilities. A uniform approximation version also gives substantial prior mass to small-network neighborhoods on the entire input cube and polynomial total sign-prediction mistakes for margin teachers on prescribed input sets.
\end{abstract}

\section{Why layerwise tails help}
A small network can approximate a complicated regression function without fitting all of its noise. We seek a Bayesian learning guarantee that competes with that approximation at every larger learner width, using one fixed dense architecture and independent weight coordinates.

Heavy-tailed fan-in priors offer many potential implementations of a smaller network. A prescribed order-one edge has probability of order $1/N$ in a width-$N$ row, but there are order $N$ possible locations for it. This compensation can make a small function neighborhood's prior probability independent of learner width.

Unbounded activations introduce another issue. A preceding hidden feature can itself be heavy-tailed. For example, with Cauchy weights and ReLU, the absolute hidden values have infinite mean; their average grows with width, and the next Cauchy readout scale can diverge. We use a strictly smaller tail index in each later layer. The next layer then requires only a fractional moment that the preceding layer possesses.

The learner does not truncate its weights or choose an architecture from a prior. Truncation appears only in a comparison measure used to lower-bound a function neighborhood's prior probability. The actual weight prior remains an independent continuous product distribution with full support.

\paragraph{Related work.}
Castillo and Egels~\cite{castillo-egels} prove adaptive posterior contraction for dense ReLU networks with independent heavy-tailed weights and deterministic architectures, under prescribed architecture and shrinkage regimes. Thus dense independent heavy-tailed priors without an architecture or sparsity hyperprior are already an established approach. Our target here is a finite-sample agnostic oracle against arbitrary smaller bounded-weight networks, with constants uniform over every prescribed finite learner width above comparator width and fan-in normalization that does not depend on sample size. The model-specific step is the width-uniform function-neighborhood calculation. The distinction is not between an ordinary and a fractional Gaussian likelihood: tempering a Gaussian likelihood is equivalent to increasing its observation variance.

\section{Model and theorem}
Fix hidden depth $H\ge1$, total depth $L=H+1$, and input dimension $d\ge1$. Let $\phi:\R\to\R$ be globally $A$-Lipschitz, with $A\ge1$ and $b=|\phi(0)|$. The activation need not be bounded or vanish at zero.

Fix a bounded output link $\psi:\R\to[-1,1]$. For the general theorem it may be any fixed Lipschitz function. For the sharper smooth-activation theorem, take either $\psi=\clip_{[-1,1]}$ or a continuously differentiable link with bounded derivative and Lipschitz derivative. Tanh is an example in both cases. The dense width-$N$ learner is
\begin{equation}
 h^1=\phi(W^1\bar x),\quad
 h^\ell=\phi(W^\ell h^{\ell-1})\ (2\le\ell\le H),\quad
 z_w=v^\top h^H,\quad f_w=\psi(z_w),\qquad \bar x=(1,x).
 \label{eq:network}
\end{equation}
The first layer includes a bias coordinate. A constant wire can compile later biases when the activation admits a nonzero constant value; this uses one additional neuron and a fixed change in coefficient bounds.

Choose tail indices satisfying
\begin{equation}
 0<\alpha_L<\cdots<\alpha_2<\min\{1,\alpha_1\},
 \qquad \alpha_1<2,\qquad a=\alpha_2.
 \label{eq:tail-indices}
\end{equation}
For each index, let
\[
 p_\alpha(t)=c_\alpha(1+t^2)^{-(1+\alpha)/2}.
\]
All prior coordinates are independent. In layer $\ell$, a coordinate has law $F^{-1/\alpha_\ell}T$, where $T$ has density $p_{\alpha_\ell}$ and $F$ is that row's fan-in: $d+1$ in the first layer and $N$ subsequently. This defines one prior $\mu_N$ for each prescribed finite architecture. The tail indices and architecture are fixed before observing data.

The joint law of $(X,Y)$ is arbitrary subject to $X\in[-1,1]^d$, a bounded conditional mean $f_0(x)=\E[Y\mid X=x]\in[-1,1]$, and conditionally centered sub-Gaussian noise $\xi=Y-f_0(X)$:
\begin{equation}
 \E[e^{t\xi}\mid X]\le e^{\sigma^2t^2/2}\quad(t\in\R).
 \label{eq:noise-assumption}
\end{equation}
Labels may be unbounded. Define the noise floor $R_0=\E\xi^2$ and squared risk $R(f)=\E(f(X)-Y)^2$. Let $S$ consist of $n\ge2$ IID observations.

For a fixed $0<\eta\le1$, use the temperature
\[
 \beta=\frac{\eta}{32(1+\sigma^2)}
\]
and the ordinary posterior
\begin{equation}
 d\Pi_{S,\beta}(w)=Z_{S,\beta}^{-1}
 e^{-\beta\sum_{i=1}^n(f_w(X_i)-Y_i)^2}\,d\mu_N(w),
 \qquad \bar f_{S,\beta}(x)=\int f_w(x)\,d\Pi_{S,\beta}(w).
 \label{eq:posterior}
\end{equation}
The Gaussian observation variance is the fixed constant $16(1+\sigma^2)/\eta$. It does not scale with $n$ or $N$. The output link is inside the likelihood.

A comparator in $\calG_{m,B}$ has the same depth, activation, first-layer bias convention and output link, width $m\le N$, and coefficients bounded by $B\ge1$. Put
\begin{align}
 D&=1+Am(B+6),& M&=(m+1)^H,\\
 U_{m,B}&=C(d+1)(m+1)^{H-1}(B+6)^{aH},&
 V_{m,B}&=C(d+1)^2D^{2H},\\
 K_{m,B}&=C\left[
 U_{m,B}V_{m,B}^{a/(2-a)}
 +(d+1)M\log\!\big(CM(d+1)(B+6)V_{m,B}\big)
 \right].
 \label{eq:resources}
\end{align}
Constants $C\ge1$ may depend on $H$, the fixed tail indices, $A,b$ and the fixed output link, but not on $d,m,B,N,n$. In the smooth case they may also depend on a Lipschitz constant for $\phi'$. Different occurrences may denote different sufficiently large constants.

\begin{theorem}[One fixed prior and likelihood for Lipschitz activations]
\label{thm:fast-oracle}
Fix a finite learner width $N$ before drawing the sample, a comparator class with $m\le N$, $B\ge1$, and $0<\delta<1$. Let
\[
 A_*=\min_{g\in\calG_{m,B}}\{R(g)-R_0\}.
\]
Write $q=a/(2-a)$. With no smoothness assumption on either activation or output link, with probability at least $1-\delta$,
\begin{equation}
 R(\bar f_{S,\beta})
 \le \int R(f_w)\,d\Pi_{S,\beta}(w)
 \le R_0+(1+\eta)A_*
 +\frac C\eta\left(
 \frac{(1+\sigma^2)\{K_{m,B}+\log(2/\delta)\}}n
 \right)^{1/(q+1)}.
 \label{eq:main-oracle}
\end{equation}
The constants are uniform over all finite $N\ge m$. The learner uses neither comparator parameters nor the unknown regression function.
If $\phi'$ is globally Lipschitz and the output link satisfies the sharper condition above, the remainder may instead be taken as
\[
 C\left(\frac{(1+\sigma^2)\{K_{m,B}+\log(2/\delta)\}}{\eta n}
 \right)^{1/(q+1)}.
\]
This improves the dependence on $\eta$; the sample-size power is the same.
\end{theorem}

\begin{theorem}[A sharp expected-risk oracle using all observations]
\label{thm:expected-oracle}
At the fixed temperature $\beta_*=1/[256(1+\sigma^2)]$, write
$q=a/(2-a)$ and $A_g=R(g)-R_0$. For every fixed comparator $g\in\calG_{m,B}$,
\begin{align}
 \E_{S_n}R(\bar f_{S_n,\beta_*})
 \le R(g)+C\bigg[
 &A_g^{q/(2q+1)}
 \left(\frac{(1+\sigma^2)K_{m,B}}{n+1}\right)^{1/(2q+1)}
 \nonumber\\
 &+\left(\frac{(1+\sigma^2)K_{m,B}}{n+1}\right)^{1/(q+1)}
 \bigg].
 \label{eq:expected-oracle}
\end{align}
If the hidden activation and output link satisfy the smooth condition of Theorem~\ref{thm:fast-oracle}, the bracket can instead be replaced by
\[
 C\left(\frac{(1+\sigma^2)K_{m,B}}{n+1}\right)^{1/(q+1)}.
\]
Both statements are uniform over all finite $N\ge m$ and concern the ordinary final posterior trained on all $n$ observations. The infimum over deterministic comparators and their resource classes may be taken directly.
\end{theorem}
\begin{proof}
The L2 comparison lemma gives risk at most
$R(g)+2\sqrt{A_g r}+r$ and KL at most $K_{m,B}r^{-q}$.
The generic all-data expectation theorem of Section~\ref{sec:sharp-expectation-proof} therefore bounds risk by
\[
 R(g)+2\sqrt{A_g r}+r+\frac{256(1+\sigma^2)K_{m,B}r^{-q}}{n+1}.
\]
Put $u=(1+\sigma^2)K_{m,B}/(n+1)$. If $u\le1$, take $r$ to be the smaller of $u^{1/(q+1)}$ and $(u/\sqrt{A_g})^{1/(q+1/2)}$, omitting the second choice when $A_g=0$. The square-root, linear and entropy terms give the two powers in~\eqref{eq:expected-oracle}, up to constants depending on the fixed tail indices. If $u>1$, bounded excess risk gives the conclusion directly. In the smooth case the direct risk comparison is $R(g)+r$, so taking $r$ of order $u^{1/(q+1)}$ proves its sharper bound.
\end{proof}

This is an expectation guarantee with coefficient exactly one on comparator risk. For example, at $a=1/2$, smooth activations give rate $n^{-3/4}$ at arbitrary approximation error. With only Lipschitzness, the remainder is of order $A_g^{1/5}n^{-3/5}+n^{-3/4}$, suppressing comparator resources. The high-probability near-sharp theorem above retains rate $n^{-3/4}$ for all Lipschitz activations.

\begin{corollary}[Automatic adaptation to the comparator]\label{cor:structural-oracle}
The same fixed-temperature learner satisfies, with probability at least $1-\delta$,
\[
 R(\bar f_{S,\beta})-R_0
 \le\inf_{\substack{1\le m\le N\\B\in\mathbb N,\ B\ge1}}
 \left\{(1+\eta)A_{m,B}
 +\frac C\eta\left(
 \frac{(1+\sigma^2)\{K_{m,B}+\log(CmB/\delta)\}}n
 \right)^{1-a/2}\right\},
\]
where $A_{m,B}=\min_{g\in\calG_{m,B}}(R(g)-R_0)$.
\end{corollary}
\begin{proof}
Allocate failure probabilities proportional to $\delta/(m^2B^2)$ in Theorem~\ref{thm:fast-oracle} and take a union bound. The prior, likelihood and prediction rule do not depend on those proof-only allocations or on the comparator class. Integer coefficient caps cover all finite real caps by rounding upward.
\end{proof}

For example, take $\alpha_1=1$, $\alpha_2=1/2$, and strictly decreasing later indices in $(1/4,1/2)$. Then the rate is $n^{-3/4}$ for ReLU, GELU, softplus, SiLU and tanh alike. These are remainders in an oracle with factor $1+\eta$ on approximation error above the noise floor. The factor on irreducible noise is exactly one.

At bounded coefficients and fixed input dimension, the displayed resource has teacher-width degree
\begin{equation}
 \max\left\{H,\ H-1+\frac{2Ha}{2-a}\right\},
 \label{eq:width-degree}
\end{equation}
up to logarithms. For $a=1/2$ this is $\max\{H,5H/3-1\}$. At fixed $m,d$, the coefficient-bound degree is at most $aH+2Ha/(2-a)$, apart from logarithms; it tends to zero as a fixed tail index $a$ is chosen smaller. The constants can depend badly on depth and on small gaps between successive tail indices. No polynomial-depth guarantee is asserted.

\begin{corollary}[Competing with polynomially growing comparator classes]
\label{cor:growing-comparators}
Fix the depth, input dimension, activation, output link, tail indices and noise proxy. Put
\[
 q=\frac a{2-a},\qquad w=H-1+2Hq,\qquad b_*=aH+2Hq.
\]
Suppose nonnegative $u,v$ satisfy
\[
 \max\{Hu,\ wu+b_*v\}<1.
\]
For any deterministic learner-width schedule $N(n)\ge\lfloor n^u\rfloor$, use
$\eta_n=1/\log(en)$ in the prescribed likelihood. With probability at least $1-\delta$ along an IID sample sequence,
\[
 R(\bar f_{S_n,\beta_n})
 \le\min_{g\in\calG_{\lfloor n^u\rfloor,n^v}}R(g)+o(1).
\]
The remainder is uniform over the allowed input-label laws and learner widths at fixed noise proxy and confidence. The prior has no sample-size-dependent change other than the prescribed architecture; the likelihood variance grows only as $\log n$.
\end{corollary}
\begin{proof}
The explicit resource is at most
$\widetilde O(n^{\max\{Hu,wu+b_*v\}})$, so its statistical remainder times $\eta_n^{-1}=\log(en)$ tends to zero. The additional approximation-factor contribution is at most $4\eta_n$, because both the comparator and conditional mean lie in $[-1,1]$. Apply Theorem~\ref{thm:fast-oracle} with failure probabilities proportional to $\delta/n^2$ and take a union bound; the additional logarithms do not change the limit.
\end{proof}

The probability statement is for each architecture fixed before the data, with constants uniform in its width. It does not supply one common event for unrestricted data-dependent choice among infinitely many widths. Exact posterior inference need not be computationally efficient.

\section{Scalar tail estimates}
\label{sec:scalar}
We first record the coordinate estimates used by the comparison construction. All constants in this section may depend on the fixed tail exponents.

\begin{lemma}[Conditioning a coordinate to be small]\label{lem:small-coordinate}
Let $W=F^{-1/\alpha}T$, where $F\ge1$, $0<\alpha<2$, and $T$ has density $p_\alpha$. For $0<u\le1$,
\begin{align}
 -\log\Pp(|W|\le u)&\le C_\alpha u^{-\alpha}/F,\qquad
 \E[W\mid |W|\le u]=0,\\
 \E[W^2\mid |W|\le u]&\le C_\alpha u^{2-\alpha}/F.
 \label{eq:small-coordinate}
\end{align}
If $\alpha<1$, also
\begin{equation}
 \E[|W|\mid |W|\le u]\le C_\alpha u^{1-\alpha}/F.
 \label{eq:small-first-moment}
\end{equation}
For $|w|\le B+5$, the unconditioned density obeys
\begin{equation}
 p_{F,\alpha}(w)\ge\frac{c_\alpha'}{F(B+6)^{1+\alpha}}.
 \label{eq:density-lower}
\end{equation}
\end{lemma}
\begin{proof}
Put $v=F^{1/\alpha}u$. For $v\le1$, the small-ball probability is at least $c_\alpha'v$, and $v^\alpha\log(1/(c_\alpha'v))$ is bounded. This proves the logarithmic estimate after using $F=v^\alpha u^{-\alpha}$. Conditional absolute moments of orders one and two are at most $u$ and $u^2$; multiplying by $F$ and using $Fu^\alpha\le1$ gives the stated estimates in this range.

For $v\ge1$, the small-ball probability is bounded below, the tail is at most $C_\alpha v^{-\alpha}$, and integration gives truncated moments bounded by $C_\alpha v^{r-\alpha}$ for $r>\alpha$. These facts imply all the estimates after rescaling. Symmetry gives the conditional mean. Finally
\[
 p_{F,\alpha}(w)=\frac{c_\alpha}{F}
       (F^{-2/\alpha}+w^2)^{-(1+\alpha)/2},
\]
which proves~\eqref{eq:density-lower}.
\end{proof}

\begin{lemma}[Fractional moments of absolute row sums]\label{lem:absolute-sum}
Suppose $0<p<\alpha<1$. For independent fan-in-$F$ coordinates $W_j$ with tail index $\alpha$ and deterministic $c_j\ge0$,
\begin{equation}
 \E\left(\sum_j c_j|W_j|\right)^p
 \le C_{\alpha,p}\left(\frac1F\sum_j c_j^\alpha\right)^{p/\alpha}.
 \label{eq:absolute-sum}
\end{equation}
\end{lemma}
\begin{proof}
Write $K=F^{-1}\sum_jc_j^\alpha$. Each summand has tail at most $C_\alpha c_j^\alpha/(Fu^\alpha)$ and an unconditional truncated first moment at level $u$ at most $C_\alpha c_j^\alpha u^{1-\alpha}/F$. Split the sum according to whether any summand exceeds $u/2$. A union bound controls that event by $C_\alpha Ku^{-\alpha}$. On its complement, Markov's inequality applied to the truncated sum gives the same bound for the event that the sum exceeds $u$. Thus
\[
 \Pp\left\{\sum_jc_j|W_j|>u\right\}
 \le\min\{1,C_\alpha Ku^{-\alpha}\}.
\]
Integrate $pu^{p-1}$ times this bound to obtain~\eqref{eq:absolute-sum}. The case $K=0$ is immediate.
\end{proof}

\begin{lemma}[Unbiased distinguished coefficients]\label{lem:distinguished}
For a teacher coefficient $c\in[-B,B]$ and $0<\epsilon\le1$, there is an absolutely continuous law $Q_c$ with
\[
 \E W=c,\quad \operatorname{Var}(W)\le17\epsilon^2,\quad
 |W-c|\le9\epsilon,\quad |W|\le B+5,\quad |W|>\epsilon
 \quad\text{almost surely}.
\]
Its density is bounded by $1/(2\epsilon)$, and hence
\[
 \KL(Q_c\|p_{F,\alpha})
 \le\log\!\frac{C_\alpha F(B+6)^{1+\alpha}}\epsilon.
\]
\end{lemma}
\begin{proof}
For $|c|>4\epsilon$, use the uniform law on $[c-\epsilon,c+\epsilon]$. Otherwise mix the uniform laws on $[3\epsilon,5\epsilon]$ and $[-5\epsilon,-3\epsilon]$ with probabilities $(1+c/(4\epsilon))/2$ and $(1-c/(4\epsilon))/2$. The latter has variance at most $(16+1/3)\epsilon^2$. The support and density bounds are immediate, and the KL estimate follows from~\eqref{eq:density-lower}.
\end{proof}

The arguments use only symmetric densities comparable above and below to $(1+t^2)^{-(1+\alpha)/2}$. The same conclusions therefore hold for any such coordinate family, with modified constants. This includes symmetric $\alpha$-stable densities, whose positivity and power-law tails give the comparison; see~\cite{zolotarev}.

\section{A conditional comparison measure}
\label{sec:comparison}
Fix a comparator $g\in\calG_{m,B}$ and precision $0<\epsilon\le1$. We construct a measure $\nu$ absolutely continuous with respect to the original independent-coordinate prior. It is independent of the data. Unlike the prior, the comparison measure can have dependencies across layers.

\subsection{Embeddings}
Unroll the comparator into a formula by duplicating a neuron whenever different parents use it. There are $m^{H-\ell+1}$ node occurrences in hidden layer $\ell$ and one scalar readout. The total number of nodes is at most $M=(m+1)^H$.

If $N\ge M$, partition each hidden layer into disjoint index groups, one per node occurrence, with every group of size at least $N/(2M)$. Choose one index uniformly and independently from each group. If $m\le N<M$, use the original comparator at fixed physical positions instead. Its number of selected rows at each layer is no larger than the corresponding formula count.

Conditional on an assignment, selected first-layer rows have independent uniform coordinates of radius $\epsilon$ around the teacher coefficients. At every selected later row, including the readout, each designated child coefficient has the law $Q_c$ from Lemma~\ref{lem:distinguished}. The other coordinates will be conditioned according to the preceding layer.

\subsection{Adaptive cutoffs}
For every neuron, selected or not, define the absolute path bound
\begin{equation}
 R_j^1=b+A\sum_{k=0}^d|W_{jk}^1|,\qquad
 R_j^\ell=b+A\sum_{k=1}^N|W_{jk}^\ell|R_k^{\ell-1},
 \qquad T_j^\ell=\max\{1,R_j^\ell\}.
 \label{eq:path-bound}
\end{equation}
Since $|\phi(t)|\le b+A|t|$, this bounds $\sup_{x\in[-1,1]^d}|h_j^\ell(x)|$.

At every undesignated coordinate entering a selected row in layer $\ell\ge2$, use its prior law conditioned on
\begin{equation}
 |W_{jk}^\ell|\le\epsilon/T_k^{\ell-1}.
 \label{eq:adaptive-cutoff}
\end{equation}
The scalar readout is a selected row and uses the same rule. Unselected rows keep their native conditional product law.

This gives a valid feedforward probability measure: all cutoffs are measurable functions of preceding layers. Conditional on preceding layers, the incoming coordinates of a row are independent. Every background coordinate is centered; every designated coefficient has its exact teacher mean and is independent of the preceding layers.

\begin{lemma}[Embedding entropy still cancels]\label{lem:embedding-entropy}
Different full index assignments have disjoint comparison supports. For their uniform mixture,
\[
 \KL(\nu\|\mu_N)
 =\text{average component KL}-\text{assignment entropy}.
\]
Each designated internal or readout edge cancels its $\log N$ density cost against its child's index entropy, leaving at most a $\log(2M)$ cost.
\end{lemma}
\begin{proof}
In a selected parent row, a designated edge has magnitude greater than $\epsilon$, while every undesignated edge has magnitude at most $\epsilon/T_k\le\epsilon$. Starting at the scalar readout, the unique large coordinate in each designated child group therefore identifies that child's chosen index. Repeating this step down the network decodes the whole assignment. The adaptive cutoffs cannot affect this decoding.

On a component's support the mixture density equals its component density times its assignment probability. Integrating the logarithm proves the entropy identity. Every non-root formula node has one parent, so its group entropy cancels exactly one designated-edge $\log N$ term. In the small-width branch there is no assignment entropy, but $\log N\le\log M$ gives the same upper bound directly.
\end{proof}

The remaining task is to bound the expected entropy of the adaptive background cutoffs. The next section establishes the moments needed for this and for the risk comparison.

\section{Path moments, entropy and risk}
\label{sec:moments}
Throughout this section, first fix an embedding. All bounds will be uniform over embeddings, so they also hold for their mixture.

\begin{lemma}[Selected first moments and native fractional moments]
\label{lem:path-moments}
Let $S_\ell$ be the largest $\E R_j^\ell$ over selected hidden neurons, and put
\[
 P_\ell=\frac1N\sum_{j=1}^N\E(T_j^\ell)^{\alpha_{\ell+1}}
 \qquad(1\le\ell\le H).
\]
Under the comparison law,
\begin{align}
 S_\ell&\le C_\ell(d+1)(B+6)D^{\ell-1},\\
 P_\ell&\le C_\ell\left[(d+1)
       +\{(d+1)(B+6)D^{\ell-1}\}^{\alpha_{\ell+1}}\right].
 \label{eq:path-moments}
\end{align}
where $C_\ell$ may depend on depth, the tail indices and the activation, but not on $d,m,B,N,\epsilon$.
\end{lemma}
\begin{proof}
A selected first-layer row satisfies $R_j^1\le b+A(d+1)(B+1)$. For an unselected row, let $p=\alpha_2<\min\{1,\alpha_1\}$. Subadditivity and the finite $p$th coordinate moment give
\[
 \E(R_j^1)^p
 \le C\{1+(d+1)\E|W_{jk}^1|^p\}
 \le C\{1+(d+1)^{1-p/\alpha_1}\}\le C(d+1).
\]
This proves the initial case, including $P_1$.

For a later selected row, conditional on preceding layers, Lemma~\ref{lem:small-coordinate} gives
\begin{align*}
 \E\sum_{\mathrm{background}}|W_{jk}^\ell|R_k^{\ell-1}
 &\le \frac{C\epsilon^{1-\alpha_\ell}}N
          \sum_k(T_k^{\ell-1})^{\alpha_\ell}.
\end{align*}
Here and below conditional expectations are with respect to the current incoming row. Designated coefficients have absolute value at most $B+5$. Therefore
\begin{equation}
 S_\ell\le b+Am(B+5)S_{\ell-1}+CA P_{\ell-1},
 \label{eq:selected-first-moment}
\end{equation}
using $\epsilon\le1$.

At an unselected row, take $p=\alpha_{\ell+1}<\alpha_\ell<1$. Lemma~\ref{lem:absolute-sum} conditionally on preceding layers gives
\[
 \E[(T_j^\ell)^p\mid\text{preceding layers}]
 \le C\left[1+
   \left(\frac1N\sum_k(T_k^{\ell-1})^{\alpha_\ell}\right)^{p/\alpha_\ell}
 \right].
\]
Jensen's inequality applies because $p/\alpha_\ell<1$. At a selected row, $\E(T_j^\ell)^p\le1+S_\ell^p$. Consequently
\begin{equation}
 P_\ell\le C\{1+P_{\ell-1}^{p/\alpha_\ell}+S_\ell^p\}.
 \label{eq:fractional-moment}
\end{equation}
Induction first bounds $S_\ell$ by the displayed geometric majorant, because all powers in the fractional recurrence are below one. To retain the sharper bound on $P_\ell$, raise its preceding bound to $p/\alpha_\ell$: the teacher term becomes $\{(d+1)(B+6)D^{\ell-2}\}^{p}$, while the dimension term is at most $d+1$. Together with $S_\ell^p$, these are bounded by~\eqref{eq:path-moments}.
\end{proof}

The strict tail decrease has a precise use here. We request a moment of order $\alpha_{\ell+1}$ from a native layer of index $\alpha_\ell$, never a moment at its own tail index. Selected nodes have stronger moments because their incoming background coordinates are conditionally truncated.

\begin{lemma}[The comparison entropy]\label{lem:comparison-entropy}
For the comparison measure, uniformly over every finite $N\ge m$,
\begin{equation}
 \KL(\nu\|\mu_N)
 \le C\left[
 U_{m,B}\epsilon^{-a}
 +(d+1)M\log\frac{CM(d+1)(B+6)}\epsilon
 \right],
 \qquad a=\alpha_2.
 \label{eq:comparison-entropy}
\end{equation}
\end{lemma}
\begin{proof}
At a background coordinate entering layer $\ell$, its conditional KL is
$-\log\Pp\{|W|\le\epsilon/T_k^{\ell-1}\}$, bounded by
$C\epsilon^{-\alpha_\ell}(T_k^{\ell-1})^{\alpha_\ell}/N$.
Thus a selected row costs at most $C\epsilon^{-\alpha_\ell}P_{\ell-1}$ after expectation.

The number of selected rows entering layer $\ell\ge2$ is at most $m^{H-\ell+1}$, with one row at $\ell=H+1$. Since $\alpha_\ell\le a$ and $\epsilon\le1$, the total background cost is bounded by
\[
 C\epsilon^{-a}\sum_{\ell=2}^{H+1}m^{H-\ell+1}
 \left[(d+1)+\{(d+1)(B+6)D^{\ell-2}\}^{\alpha_\ell}\right].
\]
Use $D\le C_A m(B+6)$, $\alpha_\ell<1$, and $d+1\ge1$. The second summand in each term is bounded by
\[
 C(d+1)m^{H-\ell+1+\alpha_\ell(\ell-2)}
       (B+6)^{\alpha_\ell(\ell-1)}
 \le C(d+1)(m+1)^{H-1}(B+6)^{aH}.
\]
The first summands satisfy the same bound. Summing the fixed number of layers gives $C U_{m,B}\epsilon^{-a}$.

Use the relative-entropy chain rule for the feedforward conditional law. Lemma~\ref{lem:embedding-entropy} cancels the designated-edge $\log N$ terms in the large-width branch, while $\log N\le\log M$ suffices in the small-width branch. There are at most $M$ designated internal/readout edges and $M(d+1)$ first-layer coordinates. Their interval density costs give the logarithmic term in~\eqref{eq:comparison-entropy}. The moment bounds show that every conditional entropy expectation used in this chain rule is finite.
\end{proof}

\begin{lemma}[Pointwise raw perturbation moments]\label{lem:raw-perturbation}
Let $z_g$ denote the comparator's raw readout. Uniformly over the input cube,
\begin{equation}
 \E_\nu|z_w(x)-z_g(x)|^2
 \le V_{m,B}\epsilon^{2-a}.
 \label{eq:raw-rms}
\end{equation}
If $\phi'$ exists and is globally Lipschitz, the constant in $V_{m,B}$ may be enlarged so that also
\begin{equation}
 \left|\E_\nu[z_w(x)-z_g(x)]\right|
 \le V_{m,B}\epsilon^{2-a}.
 \label{eq:raw-bias}
\end{equation}
\end{lemma}
\begin{proof}
Teacher hidden values satisfy
\[
 \sup_x|h_j^{*,\ell}(x)|
 \le C(d+1)(B+6)D^{\ell-1}.
\]
Fix an input and let $v_\ell$ be the largest selected-neuron second moment of the error. In the first layer,
$v_1\le A^2(d+1)\epsilon^2/3$.

At a later selected node, write the preactivation error as
\[
 \Delta=\sum_{j=1}^m c_j(h_j-h_j^*)+Z.
\]
Conditionally on preceding layers, $Z$ is centered, and its variance is bounded by
\begin{equation}
 \E[Z^2\mid\text{preceding layers}]
 \le17\epsilon^2\sum_{\mathrm{designated}}|h_j|^2
 +\frac{C\epsilon^{2-\alpha_\ell}}N
       \sum_k(T_k^{\ell-1})^{\alpha_\ell}.
 \label{eq:adaptive-variance}
\end{equation}
For a background term, this follows because its conditional variance is at most
$C(\epsilon/T_k)^{2-\alpha_\ell}/N$, while $|h_k|^2\le T_k^2$.
No second moment of an unselected hidden feature is required.

Conditional centering removes the cross term with the preceding-layer error. Since
$|h_j|^2\le2|h_j-h_j^*|^2+2|h_j^*|^2$, Lipschitzness yields
\begin{align}
 v_\ell\le A^2\bigg[&(m^2B^2+34m\epsilon^2)v_{\ell-1}\\
 &+Cm(d+1)^2(B+6)^2D^{2\ell-4}\epsilon^2
 +C P_{\ell-1}\epsilon^{2-\alpha_\ell}\bigg].
 \label{eq:variance-recursion}
\end{align}
Induction, $\epsilon^2\le\epsilon^{2-a}$, $\alpha_\ell\le a$, and Lemma~\ref{lem:path-moments} give the following bound. In the active-noise term, $m(B+6)^2\le D^2$ absorbs the apparent extra coefficient factor; the first-layer error itself has no such factor:
\[
 v_\ell\le C_\ell(d+1)^2D^{2\ell-2}\epsilon^{2-a}.
\]
The readout obeys the same recursion with activation Lipschitz constant one. Its additional factor is bounded by a constant times $D^2$, proving~\eqref{eq:raw-rms}.

For the smooth case, let $J$ be a Lipschitz constant for $\phi'$, and let $b_\ell$ be the largest absolute selected-neuron mean error. Centering of the first-layer preactivation error gives
$b_1\le J(d+1)\epsilon^2/6$. Taylor's inequality at a later teacher preactivation gives
\[
 b_\ell\le AmB\,b_{\ell-1}+\frac J2\E\Delta^2.
\]
The already established variance bounds control the second term, yielding
$b_\ell\le C_\ell(d+1)^2D^{2\ell-2}\epsilon^{2-a}$.
The readout bias is at most $mB\,b_H$, because all its coefficient perturbations and adaptive backgrounds are conditionally centered. This proves~\eqref{eq:raw-bias}. Every estimate is uniform over embeddings and hence holds for their mixture.
\end{proof}

\begin{lemma}[A data-independent mean-square comparison]\label{lem:l2-comparison}
For each $g\in\calG_{m,B}$, $0<r\le1$ and finite $N\ge m$, there is a measure $\nu_{g,r,N}\ll\mu_N$, independent of the observations, satisfying
\begin{equation}
 \sup_x\E_\nu|f_w(x)-g(x)|^2\le r,\qquad
 \KL(\nu\|\mu_N)\le K_{m,B}r^{-q},
 \qquad q=\frac a{2-a}.
 \label{eq:l2-comparison}
\end{equation}
No smoothness assumption is needed.
\end{lemma}
\begin{proof}
A Lipschitz output link gives $\E|f_w-g|^2\le L_\psi^2V_{m,B}\epsilon^{2-a}$ by Lemma~\ref{lem:raw-perturbation}. Choose $\epsilon=(r/(L_\psi^2V_{m,B}))^{1/(2-a)}$, increasing constants so $\epsilon\le1$. Substitute into~\eqref{eq:comparison-entropy}. The power term is $C U_{m,B}V_{m,B}^{a/(2-a)}r^{-q}$, and $r^q\log(1/r)\le1/(eq)$ absorbs the logarithmic term. This is precisely the resource~\eqref{eq:resources}.
\end{proof}

\begin{lemma}[A data-independent risk comparison]\label{lem:comparison}
For each $g\in\calG_{m,B}$, every $0<r\le1$ and every finite $N\ge m$, there is $\nu_{g,r,N}\ll\mu_N$, independent of the observations, such that
\begin{equation}
 \E_\nu(f_w(x)-y)^2\le(g(x)-y)^2+r
 \quad(x\in[-1,1]^d,\ |y|\le1),\qquad
 \KL(\nu\|\mu_N)\le K_{m,B}r^{-q},
 \label{eq:risk-comparison}
\end{equation}
with exponent $q=2a/(2-a)$ in general, or $q=a/(2-a)$ if $\phi\in C^{1,1}$ and the output link has the sharper regularity specified in Section~2.
\end{lemma}
\begin{proof}
For a Lipschitz output link with constant $L_\psi$, the loss
$\ell_y(z)=(\psi(z)-y)^2$ is $4L_\psi$-Lipschitz for $|y|\le1$.
Thus~\eqref{eq:raw-rms} gives risk increase at most
$4L_\psi\sqrt{V_{m,B}}\,\epsilon^{(2-a)/2}$.
Choose $\epsilon$ to make this at most $r$ and use~\eqref{eq:comparison-entropy}. Its power term becomes
\[
 C U_{m,B}V_{m,B}^{a/(2-a)}r^{-2a/(2-a)}.
\]

For the sharper case, squared loss after the output link satisfies
\[
 \ell_y(z+u)\le\ell_y(z)+d_y(z)u+C_\psi u^2,
 \qquad |d_y(z)|\le D_\psi,
\]
with fixed constants. For clipping, one may use $D_\psi=4,C_\psi=1$, because $\ell_y(z)-z^2$ is concave. For a bounded smooth link with $|\psi'|\le L_\psi$ and $\operatorname{Lip}(\psi')\le J_\psi$, take
$D_\psi=4L_\psi$ and $C_\psi=L_\psi^2+2J_\psi$.
Equations~\eqref{eq:raw-rms}--\eqref{eq:raw-bias} therefore give risk increase at most $C V_{m,B}\epsilon^{2-a}$. The chosen cutoff now makes the entropy power term
\[
 C U_{m,B}V_{m,B}^{a/(2-a)}r^{-a/(2-a)}.
\]
In both cases the same resource~\eqref{eq:resources} suffices. The remaining logarithm of $1/r$ is absorbed using $r^q\log(1/r)\le1/(eq)$. All cutoff choices are at most one after increasing the fixed constants.
\end{proof}

The pointwise comparison concerns bounded reference labels. To handle unbounded observed labels, the next section applies it with $y=f_0(x)$, then concentrates the comparison-averaged relative loss.

\section{A fixed-temperature oracle for noisy regression}
\label{sec:noisy-regression}
We now prove Theorem~\ref{thm:fast-oracle}. The central quantity is the excess risk above the regression function, which removes the irreducible noise from the complexity penalty.

Write $f_0(x)=\E[Y\mid X=x]$, $\xi=Y-f_0(X)$, and
\[
 R_0=\E\xi^2,\qquad
 A(f)=\E(f(X)-f_0(X))^2=R(f)-R_0.
\]
The assumption is $|f_0|\le1$ and
\begin{equation}
 \E[e^{t\xi}\mid X]\le e^{\sigma^2t^2/2}
 \quad\text{for every }t\in\R.
 \label{eq:conditional-subgaussian}
\end{equation}
No Gaussian assumption is needed. Bounded labels in $[-1,1]$ satisfy this with $\sigma^2=1$ by conditional Hoeffding. Gaussian label noise of variance at most $\sigma^2$ is another example.

Let $s=1+\sigma^2$ and $\beta=\eta/(32s)$, where $0<\eta\le1$. Every network output $f_w$ remains bounded, although the labels need not be. The posterior is still~\eqref{eq:posterior}; its Gaussian observation variance is the constant $16s/\eta$, independent of sample size and learner width.

\begin{lemma}[A relative-loss PAC--Bayes inequality]\label{lem:relative-pac}
With probability at least $1-\delta/2$, simultaneously for every probability measure $\rho$ on weights,
\begin{equation}
 c_\beta^{\mathrm{rel}}\E_\rho A(f_w)
 \le\E_\rho[\whR_S(f_w)-\whR_S(f_0)]
       +\frac{\KL(\rho\|\mu_N)+\log(2/\delta)}{\beta n},
 \label{eq:relative-pac}
\end{equation}
where
\[
 c_\beta^{\mathrm{rel}}=\frac{1-e^{-4\beta(1-2\beta\sigma^2)}}{4\beta}
       \ge1-4s\beta=1-\eta/8.
\]
\end{lemma}
\begin{proof}
Fix $f=f_w$ and put $e=f-f_0$. The difference of squared losses is $e(X)^2-2\xi e(X)$. Conditional sub-Gaussianity gives
\[
 \E[e^{-\beta\{(f(X)-Y)^2-(f_0(X)-Y)^2\}}\mid X]
 \le e^{-\beta(1-2\beta\sigma^2)e(X)^2}.
\]
Since $0\le e^2\le4$, convexity on that interval yields
\[
 \E e^{-\beta\{(f(X)-Y)^2-(f_0(X)-Y)^2\}}
 \le e^{-\beta c_\beta^{\mathrm{rel}} A(f)}.
\]
The lower bound on $c_\beta^{\mathrm{rel}}$ follows from $1-e^{-u}\ge u-u^2/2$ and the chosen range of $\beta$. Independence, integration against the prior, Markov's inequality and the entropy variational inequality prove~\eqref{eq:relative-pac}. The individual squared losses have finite expectation under~\eqref{eq:conditional-subgaussian}; their difference has the exponential moment just established.
\end{proof}

The comparison measure may itself be averaged before applying concentration. This is useful because the perturbation lemma is naturally a statement about its expected risk.

\begin{lemma}[Concentration for a fixed comparison law]\label{lem:comparison-concentration}
Let $\nu$ be any fixed probability measure on bounded predictions, independent of the data, and let $a=\E_\nu A(f_w)$. With probability at least $1-e^{-\ell}$,
\begin{equation}
 \E_\nu[\whR_S(f_w)-\whR_S(f_0)]
 \le a+C\sqrt{sa\ell/n}+Cs\ell/n.
 \label{eq:comparison-bernstein}
\end{equation}
\end{lemma}
\begin{proof}
Set
\[
 v(x)=\E_\nu(f_w(x)-f_0(x))^2,\qquad
 b(x)=\E_\nu(f_w(x)-f_0(x)).
\]
Then $0\le v\le4$, $b^2\le v$, and one comparison-averaged excess loss is $L=v(X)-2\xi b(X)$, of mean $a$. For $0\le t\le1/(16s)$, conditional sub-Gaussianity gives
\[
 \E e^{t(L-a)}
 \le e^{-ta}\E e^{(t+2\sigma^2t^2)v(X)}
 \le\exp\left\{a\left[\frac{e^{4(t+2\sigma^2t^2)}-1}{4}-t\right]\right\}
 \le e^{6sa t^2}.
\]
For the last inequality, put $u=t+2\sigma^2t^2\le9t/8$, use $4u\le1$ and $e^{4u}-1\le4u+16u^2$. Chernoff's bound for the independent observations then gives~\eqref{eq:comparison-bernstein}, with an absolute constant. Equivalently, optimize over $t\le1/(16s)$ to obtain a square-root term and a linear term.
\end{proof}

\begin{proof}[Proof of Theorem~\ref{thm:fast-oracle}]
Fix a deterministic risk-minimizing $g\in\calG_{m,B}$ and write $A_*=A(g)$. The minimum exists by compactness of the bounded coefficient cube and dominated convergence. For $0<r\le1$, Lemma~\ref{lem:l2-comparison} supplies a data-independent law with
\[
 \E_\nu\E_X|f_w(X)-g(X)|^2\le r,\qquad
 \KL(\nu\|\mu_N)\le K_{m,B}r^{-q},\quad q=a/(2-a).
\]
The squared triangle inequality, with its freely chosen approximation slack, gives
\begin{equation}
 a_\nu:=\E_\nu A(f_w)
 \le(1+\eta/8)A_*+(1+8/\eta)r.
 \label{eq:l2-triangle}
\end{equation}
This is the step that removes any smoothness requirement.

At the Gibbs posterior, the right side of~\eqref{eq:relative-pac} equals
$F_{S,\beta}-\whR_S(f_0)+\log(2/\delta)/(\beta n)$.
The free-energy variational formula and the comparison law bound it by
\[
 \E_\nu[\whR_S(f_w)-\whR_S(f_0)]
       +\frac{K_{m,B}r^{-q}+\log(2/\delta)}{\beta n}.
\]
Apply Lemma~\ref{lem:comparison-concentration} with $\ell=\log(2/\delta)$ and Young's inequality. On the intersection of the two events,
\[
 c_\beta^{\mathrm{rel}}\E_{\Pi_{S,\beta}}A(f_w)
 \le(1+\eta/8)a_\nu
       +\frac{Cs\ell}{\eta n}
       +\frac{32sK_{m,B}r^{-q}}{\eta n}.
\]
Substitute~\eqref{eq:l2-triangle} and divide by
$c_\beta^{\mathrm{rel}}\ge1-\eta/8$. For $0<\eta\le1$,
\[
 \frac{(1+\eta/8)^2}{1-\eta/8}\le1+\eta.
\]
All other terms are bounded by
\[
 \frac C\eta\left(r+\frac{sK_{m,B}r^{-q}}n+\frac{s\ell}n\right).
\]
Let $u=[s\{K_{m,B}+\ell\}/n]^{1/(q+1)}$. If $u\le1$, take $r=u$, obtaining remainder $Cu/\eta$. If $u>1$, the bound $A(f_w)\le4$ gives the result after increasing $C$. Adding $R_0$ and applying Jensen proves the general theorem.

For the optional smooth refinement, Lemma~\ref{lem:comparison} instead gives
$a_\nu\le A_*+r$ with the same exponent $q=a/(2-a)$. The same calculation then leaves
$C[r+s(K_{m,B}r^{-q}+\ell)/(\eta n)]$. Choosing
$r=[s\{K_{m,B}+\ell\}/(\eta n)]^{1/(q+1)}$ proves the sharper dependence on $\eta$, with the bounded-risk fallback when this exceeds one.
\end{proof}

The unknown regression function appears only inside the proof. The learner uses neither $f_0$, $R_0$, nor the comparator. The prescribed likelihood variance depends only on the chosen approximation factor $1+\eta$ and a known conditional sub-Gaussian proxy $\sigma^2$; an upper bound on that proxy is sufficient.

% Independent proof module, 18 November 2026.
% This file is a review artifact, not an edit to an owner manuscript.
% The leave-one-out strategy is motivated by Hogsgaard--Rebeschini--Wegel,
% arXiv:2607.02247 (2026). Their stated theorem is for finite dictionaries
% and bounded losses. The arbitrary-prior and conditional sub-Gaussian
% extension below is derived here; it should receive independent checking.

\section{A sharp expectation oracle for the ordinary final posterior}
\label{sec:sharp-expectation-proof}

Let $\pi$ be any probability prior on jointly measurable functions
$f_w:\mathcal X\to[-1,1]$. Suppose the observations are IID,
$f_0(x)=\mathbb E[Y\mid X=x]\in[-1,1]$, and
$\xi=Y-f_0(X)$ is conditionally sub-Gaussian with proxy $\sigma^2$.
Set $s=1+\sigma^2$ and fix $T\ge256s$.
For a sample $S_n$ of size $n$, let
\[
 dQ_{S_n}(w)=Z_{S_n}^{-1}
   \exp\left\{-\frac1T\sum_{i=1}^n(f_w(X_i)-Y_i)^2\right\}d\pi(w),
 \qquad \bar f_{S_n}=\int f_w\,dQ_{S_n}(w).
\]
This is the ordinary posterior with Gaussian observation variance $T/2$,
using all $n$ observations. No averaging over sample sizes is performed.

\begin{theorem}[Sharp all-data expectation bound]\label{thm:generic-expectation}
For every deterministic comparison law $\rho\ll\pi$,
\[
 \mathbb E_{S_n} R(\bar f_{S_n})
 \le \int R(f_w)\,d\rho(w)
        +\frac{T\operatorname{KL}(\rho\|\pi)}{n+1}.
\]
Consequently the infimum over all such comparison laws may be taken on
the right-hand side. The theorem requires neither convexity of the
support of $\pi$ nor a Bernstein condition relative to the comparator.
\end{theorem}

\begin{proof}
We first prove an averaged one-observation inequality. Fix a probability
law $q$ on $U\in[-1,1]$, with mean $\mu$ and variance $v$. For a real
label $y$, let $L_y(u)=(u-y)^2$ and
\[
 dp_y(u)=\frac{e^{-L_y(u)/T}}{\int e^{-L_y/T}\,dq}\,dq(u).
\]
For two independent draws $U,U'$ from $q$, the covariance identity gives
\begin{align*}
 \int L_y\,dq-\int L_y\,dp_y
 &=\frac{\mathbb E[(L_y(U)-L_y(U'))
         (e^{-L_y(U')/T}-e^{-L_y(U)/T})]}{2\int e^{-L_y/T}\,dq}.
\end{align*}
Since
\[
 |L_y(u)-L_y(u')|\le2(1+|y|)|u-u'|,
 \qquad \operatorname{osc}_{[-1,1]}L_y\le4(1+|y|),
\]
the mean-value theorem bounds the last display by
\[
 \frac{4(1+|y|)^2}{T}
       e^{4(1+|y|)/T}v.
\]
Also $\int L_y\,dq=L_y(\mu)+v$. Hence
\begin{equation}
 \int L_y\,dp_y\ge L_y(\mu)+
 \left[1-\frac{4(1+|y|)^2}{T}e^{4(1+|y|)/T}\right]v.
 \label{eq:one-step-pointwise}
\end{equation}

Now let $Y=a+\xi$, where $|a|\le1$ and $\xi$ is centered
$\sigma^2$-sub-Gaussian, independently of the fixed law $q$.
The elementary estimates
\[
 \mathbb E e^{t|\xi|}\le2e^{\sigma^2t^2/2},
 \qquad \mathbb E\xi^4\le16\sigma^4
\]
and Cauchy--Schwarz imply
\begin{align*}
 &\mathbb E\left[(1+|Y|)^2e^{4(1+|Y|)/T}\right]\\
 &\quad\le e^{8/T}
       \mathbb E[(8+2\xi^2)e^{4|\xi|/T}]\\
 &\quad\le e^{8/T}
       \left[16e^{8\sigma^2/T^2}
              +8\sqrt2\,\sigma^2e^{16\sigma^2/T^2}\right]
 \le32s
 \qquad(T\ge256s).
\end{align*}
Thus the expected coefficient subtracted in
\eqref{eq:one-step-pointwise} is at most $128s/T\le1/2$.
After averaging over $Y$ we obtain
\begin{equation}
 \mathbb E_Y L_Y(\mu)
 \le\mathbb E_Y\int L_Y\,dp_Y-\tfrac12v
 \le\mathbb E_Y\int L_Y\,dp_Y.
 \label{eq:averaged-one-step}
\end{equation}
The same statement holds conditionally on any input $X=x$ and on an
independent training sample.

Introduce $n+1$ IID observations, retaining the same fixed architecture,
prior and $T$ used by the $n$-observation learner. Let $Q^+$ be their full
posterior and $Q^{-i}$ the posterior after removing observation $i$.
Conditional on the other observations and $X_i$, push $Q^{-i}$ forward
through $w\mapsto f_w(X_i)$ to obtain $q$. It is independent of $Y_i$
given this conditioning. Updating $q$ by the omitted squared-loss factor
gives exactly the corresponding pushforward of $Q^+$.
Applying \eqref{eq:averaged-one-step} and then averaging yields
\[
 \mathbb E (\bar f^{-i}(X_i)-Y_i)^2
 \le\mathbb E\int(f_w(X_i)-Y_i)^2\,dQ^+(w).
\]
Sum over $i=1,\ldots,n+1$. By exchangeability and independence, the
left-hand average is $\mathbb E_{S_n}R(\bar f_{S_n})$.
For each realized enlarged sample the Gibbs variational inequality gives
\begin{align*}
 &\int\sum_{i=1}^{n+1}(f_w(X_i)-Y_i)^2\,dQ^+(w)\\
 &\quad\le\int\sum_{i=1}^{n+1}(f_w(X_i)-Y_i)^2\,d\rho(w)
             +T\operatorname{KL}(\rho\|\pi),
\end{align*}
where the nonnegative term $T\operatorname{KL}(Q^+\|\pi)$ was dropped.
Take expectations and divide by $n+1$ to prove the theorem.
All integrals are finite: the outputs are bounded and the labels have
finite second moment. The posterior normalizers are positive, and the
loss oscillation over parameters is finite for every realized sample.
\end{proof}

\paragraph{Attribution.}
The leave-one-out strategy and a pointwise inverse-tilting inequality
for bounded losses are from M.~M.~H\o gsgaard, P.~Rebeschini and
T.~Wegel, \emph{Aggregation with Exponential Weights is Optimal in
Expectation}, arXiv:2607.02247, July 2026~\cite{hogsgaard}. Their main statement uses a
finite uniform dictionary and bounded losses. Above, sums are replaced
by integrals and the variational comparison is written in KL form;
the pointwise bounded-loss step is replaced by an averaged cooling
inequality for bounded predictions and conditionally sub-Gaussian
labels. The arbitrary-prior and conditional sub-Gaussian extension is proved here.

\section{A coefficient-one oracle for bounded labels}
\label{sec:calibrated}
The fixed-temperature theorem separates the noise floor and allows the factor $1+\eta$ on approximation error. For bounded labels, the same prior also gives a coefficient-one oracle through an observable temperature calibration. No smoothness is needed, but its general-risk rate is slower.

Assume throughout this section that $Y\in[-1,1]$. Define
\[
 F_S(\beta)=-\frac1{\beta n}\log Z_{S,\beta},\quad
 \mathcal B_n=\{2^{-j}:3\le j\le\lceil\log_2(64n)\rceil\},\quad
 \ell=\log(2|\mathcal B_n|/\delta),\quad
 c_\beta=\frac{1-e^{-4\beta}}{4\beta}.
\]
Choose a grid point minimizing the observable score
\begin{equation}
 \mathcal C_S(\beta)=\frac{F_S(\beta)+\ell/(\beta n)}{c_\beta}.
 \label{eq:calibrated-score}
\end{equation}
Break ties by the grid order. This selects one Gaussian likelihood and one ordinary posterior; it does not average across temperatures.

\begin{theorem}[An observable coefficient-one oracle]\label{thm:coefficient-one}
Let $q=a/(2-a)$, $T=K_{m,B}+\ell$, and $u=(T/n)^{1/(q+1)}$.
For every fixed architecture and comparator class as in Theorem~\ref{thm:fast-oracle}, with probability at least $1-\delta$,
\begin{equation}
 R(\bar f_{S,\widehat\beta})
 \le \int R(f_w)\,d\Pi_{S,\widehat\beta}(w)
 \le R_*+C\{\sqrt{R_*u}+u\},
 \qquad R_*=\min_{g\in\calG_{m,B}}R(g).
 \label{eq:coefficient-one}
\end{equation}
The rule uses neither the comparator nor $K_{m,B}$ or $R_*$. Constants are uniform over learner width.
\end{theorem}
\begin{proof}
For $0\le l\le4$, convexity gives
$e^{-\beta l}\le1-(1-e^{-4\beta})l/4$.
For every fixed network, independence therefore implies
\[
 \E_S\exp\{\beta n[c_\beta R(f_w)-\whR_S(f_w)]\}\le1.
\]
Integrate against the prior, use Markov's inequality and relative entropy, and union-bound over the grid. With probability at least $1-\delta/2$, simultaneously at all grid points and for all posterior laws,
\[
 c_\beta\E_\rho R(f_w)\le\E_\rho\whR_S(f_w)
       +\frac{\KL(\rho\|\mu_N)+\ell}{\beta n}.
\]
At the Gibbs posterior this gives $\E_{\Pi_{S,\beta}}R(f_w)\le\mathcal C_S(\beta)$.

Fix a population-risk-minimizing $g$ and write $R=R(g)$. If $u>1$, bounded risk gives the theorem after increasing $C$. Suppose $u\le1$. Lemma~\ref{lem:l2-comparison} at mean-square precision $u$ gives, for every training sample,
\[
 \E_\nu\whR_S(f_w)
 \le\whR_S(g)+2\sqrt{\whR_S(g)u}+u,
 \qquad \KL(\nu\|\mu_N)\le K_{m,B}u^{-q}.
\]
This is just Cauchy--Schwarz applied to the cross term of the squared loss.

Take the analytical temperature
\[
 \beta_*=\frac1{16}\sqrt{\frac u{R+u}}.
\]
It is at most $1/16$ and at least $1/(36n)$, because $R\le4$ and $u\ge1/n$. The grid contains a $\beta\in[\beta_*/2,\beta_*]$.
The free-energy variational formula gives
\[
 F_S(\beta)+\frac\ell{\beta n}
 \le\whR_S(g)+2\sqrt{\whR_S(g)u}+u+\frac{T u^{-q}}{\beta n}.
\]
Since $u^{q+1}=T/n$, the last term equals $u/\beta$ and is at most
$32\sqrt{u(R+u)}\le32(\sqrt{Ru}+u)$.
Also $\beta R\le\sqrt{Ru}/16$ and $c_\beta^{-1}\le1+3\beta$.

For the one fixed teacher, bounded-loss Bernstein concentration gives, on another event of probability at least $1-\delta/2$,
\[
 \sqrt{\whR_S(g)}\le\sqrt R+C\sqrt{\ell/n}.
\]
Because $\ell/n\le u$, substituting this bound in the displayed score leaves at most $R+C(\sqrt{Ru}+u)$. The selected score is no larger. Jensen gives the ordinary posterior-mean claim.
\end{proof}

When $R_*=0$, the coefficient-one oracle has the same $n^{-(1-a/2)}$ rate as the fixed-temperature theorem. For positive comparator risk, its displayed rate is $n^{-(1-a/2)/2}$, with a square-root dependence on that risk. The fixed-temperature theorem can be substantially better when most of the test loss is irreducible label noise.

\section{Uniform function neighborhoods and total sign mistakes}
\label{sec:uniform}
The same comparison construction approximates a comparator uniformly on the entire input cube. Here the absolute first moment of truncated background coordinates is more useful than their centered variance.

\begin{lemma}[Expected uniform approximation]\label{lem:uniform-approximation}
For a constant of the same allowed dependence as above, put
\[
 W_{m,B}=C(d+1)D^H.
\]
The comparison law at precision $\epsilon$ satisfies
\begin{equation}
 \E_\nu\sup_{x\in[-1,1]^d}|z_w(x)-z_g(x)|
 \le W_{m,B}\epsilon^{1-a}.
 \label{eq:uniform-approximation}
\end{equation}
\end{lemma}
\begin{proof}
First fix an embedding. Let $e_\ell$ be the largest expected uniform hidden-coordinate error at a selected node. The first-layer error is deterministically at most $A(d+1)\epsilon$. The triangle inequality, $|W-c|\le9\epsilon$ on designated edges, and the adaptive truncated first-moment estimate give
\[
 e_\ell\le AmB\,e_{\ell-1}
      +9Am\epsilon S_{\ell-1}
      +CA\epsilon^{1-\alpha_\ell}P_{\ell-1}.
\]
Because $\alpha_\ell\le a<1$, both $\epsilon$ and $\epsilon^{1-\alpha_\ell}$ are at most $\epsilon^{1-a}$. The moment bounds and $D\ge Am(B+6)$ yield
\[
 e_\ell\le C_\ell(d+1)D^{\ell-1}\epsilon^{1-a}.
\]
The coefficient factor in the active error is absorbed by $m(B+6)\le D$. The readout obeys the same triangle inequality with activation constant one, giving~\eqref{eq:uniform-approximation}. Averaging over embeddings preserves the bound.
\end{proof}

\begin{theorem}[Width-uniform mass of a uniform function neighborhood]
\label{thm:uniform-mass}
For every $0<t\le1$ and every finite $N\ge m$,
\begin{equation}
 \mu_N\left\{\sup_{[-1,1]^d}|z_w-z_g|\le t\right\}
 \ge e^{-Q_{m,B,t}},
 \label{eq:uniform-mass}
\end{equation}
where one may take
\begin{equation}
 Q_{m,B,t}
 =C\left[
 U_{m,B}\left(\frac{2W_{m,B}}t\right)^{a/(1-a)}
 +(d+1)M\log\!\left(
 \frac{CM(d+1)(B+6)W_{m,B}}t\right)\right].
 \label{eq:uniform-mass-resource}
\end{equation}
This resource is polynomial in $d,m,B,1/t$ at fixed depth and tail indices, without a dependence on the number of input points.
\end{theorem}
\begin{proof}
Choose $\epsilon=(t/(2W_{m,B}))^{1/(1-a)}\le1$. Markov's inequality and Lemma~\ref{lem:uniform-approximation} give comparison probability at least $1/2$ to the uniform approximation event. Its comparison KL is bounded by Lemma~\ref{lem:comparison-entropy}. For an event $E$ with $\nu(E)\ge1/2$, binary relative entropy gives
\[
 \KL(\nu\|\mu_N)\ge\frac12\log(1/\mu_N(E))-\log2.
\]
Substitution proves~\eqref{eq:uniform-mass-resource}, after changing constants.
\end{proof}

\begin{corollary}[A margin function receives substantial sign prior mass]
\label{cor:sign-mass}
Let $\mathcal X$ be a nonempty compact subset of the input cube, and suppose
$f:\mathcal X\to\{\pm1\}$ has a comparator satisfying
$f(x)z_g(x)\ge2\kappa$ on $\mathcal X$, with $\kappa>0$.
Use $t=\min\{1,\kappa/2\}$ in~\eqref{eq:uniform-mass-resource}. Then
\[
 \mu_N\{\operatorname{sign}z_w=f\text{ on }\mathcal X,\ 
                      \min_{\mathcal X}|z_w|>\kappa\}
 \ge e^{-Q_{m,B,t}}.
\]
The same lower bound holds for the sign-function prior if the weight prior is initially conditioned on $\min_{\mathcal X}|z_w|>\kappa$.
\end{corollary}
\begin{proof}
On the uniform approximation event, $f(x)z_w(x)\ge2\kappa-t>\kappa$.
The conditioning event has positive prior mass, and its probability is at most one. Compactness and continuity make these uniform events measurable; the target sign is continuous on $\mathcal X$ because it has a positive teacher margin.
\end{proof}

The set $\mathcal X$ may be the Boolean cube, a disconnected continuous support, or another compact domain. A continuous nonconstant sign function cannot have a nonzero margin everywhere on a connected domain; the corollary does not remove that elementary restriction.

\begin{corollary}[Every-sequence mistake bounds]\label{cor:online}
Use either sign prior in Corollary~\ref{cor:sign-mass}, and condition exactly on each noiseless sign observation $f(X_t)$. On every realized query sequence,
\[
 \sum_{t\ge1}-\log p_t\le Q_{m,B,t_0},
 \qquad t_0=\min\{1,\kappa/2\},
\]
where $p_t$ is the posterior probability of the correct sign before its reveal. Thus majority prediction makes at most $Q_{m,B,t_0}/\log2$ mistakes, and the number of steps with $p_t\le c$ is at most $Q_{m,B,t_0}/(-\log c)$ for every $0<c<1$.

Fresh posterior-sampled predictions have expected total mistake count at most $Q_{m,B,t_0}$, provided the next query is chosen before that prediction's fresh randomness. The query may depend on all earlier predictions and observations. No IID or stationarity assumption is needed.
\end{corollary}
\begin{proof}
If $\pi$ is the induced function prior, posterior normalizers telescope along every finite realized history:
\[
 \prod_{j=1}^T p_j
 =\pi\{h:h(X_j)=f(X_j)\text{ for every }j\le T\}
 \ge\pi(f)\ge e^{-Q_{m,B,t_0}}.
\]
Take logarithms and increase $T$. A majority mistake has $p_t\le1/2$; a confidence failure has $p_t\le c$. For fresh posterior sampling, the conditional mistake probability is $1-p_t\le-\log p_t$. The tower property and monotone convergence prove the expected-count bound.
\end{proof}

This sign-learning corollary uses the same weight prior but a different observation model from the noisy-regression theorem. It concerns exact noiseless sign conditioning. The bounded output links tanh and clipping preserve the raw readout's sign.


\section{Why the output is bounded}
\label{sec:raw-mean}
A bounded output link is substantive for these priors. Without it, an ordinary posterior mean need not exist even after Gaussian-likelihood training.

\begin{proposition}[An unobserved ReLU feature can leave the raw mean undefined]
\label{prop:undefined-raw-mean}
Consider one ReLU hidden layer, a raw linear output, and a Gaussian likelihood on that raw output. Suppose the first-layer coordinate prior has an everywhere positive continuous density, and an independent readout coordinate has a symmetric power tail of index at most one. For any finite set of training inputs on the unit sphere and any different test point on that sphere, both positive and negative first moments of the raw posterior prediction at that test point are infinite. In particular, for a nonatomic sphere input law in dimension $d\ge2$, the raw posterior mean is undefined at almost every test input after every finite training sample.
\end{proposition}
\begin{proof}
Fix a test point $x$ distinct from all training points $x_i$, and put
$r=\max_i x^\top x_i<1$. The hidden weight vector $w=x$ and bias
$b_0=-(1+r)/2$ give
\[
 w^\top x_i+b_0<0\quad\text{for every training point},\qquad
 w^\top x+b_0=(1-r)/2>0.
\]
These strict inequalities persist on an open neighborhood of $(w,b_0)$.
Thus a chosen hidden neuron is exactly inactive on the training sample and positive at the test point on an event of positive prior probability. The Gaussian likelihood is positive at every finite parameter vector, so this event also has positive posterior probability.

Conditional on the other parameters on that event, the likelihood does not depend on the chosen readout coefficient at all. That coefficient retains its original symmetric heavy-tailed law. The test output is a positive multiple of it plus a finite offset. A symmetric power tail of index at most one has infinite positive and negative first moments, and the same is true after this affine transformation. Tonelli's theorem proves both infinite posterior moments. A nonatomic sphere law gives probability one to a test point distinct from a finite training set.
\end{proof}

\section{Consistency under arbitrary predetermined width growth}
\begin{corollary}[Risk consistency]\label{cor:consistency}
Assume the union of the finite comparator classes is dense in the bounded regression functions under $L^2(D_X)$. Fix the depth, tail indices, noise proxy and $0<\eta\le1$, and let $N(n)\to\infty$ be any deterministic sequence of finite learner widths. The ordinary final posterior mean at the fixed temperature of Theorem~\ref{thm:fast-oracle} satisfies
\[
 R(\bar f_{S_n,\beta})\longrightarrow R_0
 \quad\text{almost surely}
\]
along an IID sample sequence.
\end{corollary}
\begin{proof}
Choose a deterministic countable sequence of finite comparators with approximation errors tending to zero. For any one comparator, eventually $N(n)$ exceeds its width. Apply Theorem~\ref{thm:fast-oracle} with failure probability $n^{-2}$ and use Borel--Cantelli. Its remainder tends to zero, independently of the growth rate of $N(n)$. Intersect these probability-one events over the countable comparator sequence and then send its approximation error to zero. Every predictor has risk at least $R_0$, so the limit follows.
\end{proof}

For ReLU hidden layers and a final tanh or clip, the density assumption holds on compact input domains: the standard one-hidden-layer universal approximation theorem supplies continuous approximants, positive ReLU hidden features can be copied through additional fixed-depth layers, and continuous functions are dense in $L^2$ of a finite Borel measure. For a tanh output link, first shrink a continuous target slightly into $(-1,1)$ and approximate its inverse-tanh transform. The standard neural approximation fact used here is documented in~\cite{leshno}. This is a qualitative consistency corollary; the earlier explicit oracle is the quantitative statement.

\section{Scope}
The main result covers bounded and unbounded hidden activations through one argument. Its bounded output link can be smooth, including tanh; clipping is optional. It does not assert finite untruncated predictive moments for the heavy-tailed raw readout.

The tail schedule is a substantive prior choice. Smaller exponents mean heavier tails. Strict decrease is what supplies the finite fractional moments used by the proof. A same-index Cauchy prior with an unbounded ReLU hidden layer does not satisfy those moment conditions and can fail even on a constant noiseless target at a fixed likelihood temperature.

The assumptions permit arbitrary conditional sub-Gaussian label noise and model misspecification. A comparator merely approximates the conditional mean; it need not interpolate the observed labels. Neither the comparison distribution, its adaptive cutoffs nor its averaging over embeddings changes the learner's independent product prior.

The proof has received independent AI-agent mathematical review. It has not been formalized in Lean. The accompanying verification directory records the exact reviewed sources and build.
\begin{thebibliography}{9}
\bibitem{leshno} M. Leshno, V. Ya. Lin, A. Pinkus, and S. Schocken. Multilayer feedforward networks with a nonpolynomial activation function can approximate any function. \emph{Neural Networks} 6(6), 861--867, 1993.
\bibitem{hogsgaard} M. M. H\o gsgaard, P. Rebeschini, and T. Wegel. Aggregation with Exponential Weights is Optimal in Expectation. arXiv:2607.02247, 2026.
\bibitem{castillo-egels} I. Castillo and P. Egels. Posterior and Variational Inference for Deep Neural Networks with Heavy-Tailed Weights. \emph{Journal of Machine Learning Research} 26(122), 1--58, 2025.
\bibitem{zolotarev} V. M. Zolotarev. \emph{One-dimensional Stable Distributions}. Translations of Mathematical Monographs, vol. 65. American Mathematical Society, 1986.
\end{thebibliography}
\end{document}
